\documentclass{article}
\usepackage{amsmath}
\usepackage{amssymb}

\title{Ultimate LaTeX Converter Stress Test}
\author{Gemini Spark}
\date{2026-08-04}

\begin{document}

\maketitle

\section{Fixes Verification: Inline Multiplication and Cases Delimiting}
Here is a $3 \times 3$ matrix $A \in \mathbb{R}^{3 \times 3}$ and inline vector $\vec{v} \in \mathbb{C}^{3}$. Note that $3 \times 3$ uses explicit $\times$ notation.

The piecewise function with a single left curly brace and NO right closing parenthesis:
\begin{equation}
f(x, y) = \begin{cases} \frac{\sin{\left( x \cdot y \right)}}{x \cdot y}, & \text{if } x \cdot y \neq 0 \\ 1, & \text{if } x \cdot y = 0 \end{cases}
\end{equation}

\section{Operators Without Limits}
Indefinite integrals, double integrals, indefinite summations, and infinite products without explicit limits:

Indefinite single and double integrals:
\begin{equation}
\int{f(x) dx} + \iint{g(x, y) dx dy} = C
\end{equation}

Indefinite summation and infinite product without upper/lower limit blocks:
\begin{equation}
\sum{a_{n} \cdot x^{n}} = \prod{\left( 1 + b_{k} \right)}
\end{equation}

\section{Hyperbolic, Inverse, and Extremum Functions}
Hyperbolic, inverse, and logarithmic function operators:
\begin{equation}
\sinh{x} + \cosh{x} = \exp{x}
\end{equation}
\begin{equation}
\arctan{x} + \arcsin{y} = \arccos{z}
\end{equation}

Extremum operators with subscript domains:
\begin{equation}
S = \sup_{x \in \mathbb{R}}{\{ f(x) \}} + \min_{y \in [a, b]}{\{ g(y) \}} + \lim_{n \to \infty}{a_{n}}
\end{equation}

\section{Special Fonts, Vectors, and Diacritics}
Blackboard bold ($\mathbb{R}, \mathbb{C}, \mathbb{Z}, \mathbb{N}$), Calligraphic ($\mathcal{L}, \mathcal{F}$), and Fraktur ($\mathfrak{g}$) sets:
\begin{equation}
\mathcal{L}\{ f(t) \} = \int_{0}^{\infty}{e^{-s \cdot t} \cdot f(t) dt}, \quad \text{where } s \in \mathbb{C}
\end{equation}

Bold vectors ($\mathbf{v}, \bm{A}$) and diacritic accents ($\vec{r}, \dot{x}, \ddot{y}, \hat{\theta}, \bar{z}, \tilde{f}$):
\begin{equation}
\vec{r} = \mathbf{x} + \mathbf{y}, \quad \ddot{x} + 2\kappa \cdot \dot{x} + \omega_{0}^{2} \cdot x = \hat{f}(t)
\end{equation}

\end{document}
